% CAPITOLO 6 — PERSISTENCY
\chapter{Persistenza}
\label{cap:persistenza}

La persistenza è intercambiabile a runtime (NFR-01): tre backend mutuamente
esclusivi sono selezionati tramite il pattern \textbf{Abstract Factory}
(\texttt{DAOFactory} + una concrete factory per modalità).

\section{Layer di persistenza DB (MySQL/JDBC)}
Il livello di persistenza DB utilizza \textbf{MySQL} come DBMS relazionale. La
gestione dei dati è affidata a classi DAO specifiche (es.
\texttt{JDBCUserDAO}, \texttt{JDBCOrdineDAO}), che eseguono query SQL per le
operazioni CRUD sulle entità.

\textbf{Entità principali:}
\begin{itemize}
    \item \textbf{Utenti}: registrati, con i ruoli (Cliente e Venditore);
    \item \textbf{Ordini}: record transazionali con stato, totale e voci;
    \item \textbf{Prodotti}: articoli del marketplace (catalogo e inventario);
    \item \textbf{Buoni}: codici sconto (percentuale o importo fisso).
\end{itemize}

Le transazioni sono atomiche (\texttt{ConnectionManager} usa
\texttt{setAutoCommit(false)} con \texttt{commit/rollback} per la riduzione
scorte).

\textbf{Configurazione}: il sistema è configurato via il file
\texttt{config.properties} impostando \texttt{persistence\_type = DB} e
fornendo le credenziali MySQL necessarie.

\section{Layer di persistenza FS (File System)}
Il layer di persistenza FS memorizza i dati in file \textbf{JSON} nella
cartella \texttt{data/}. Ogni entità è mappata su un file dedicato
(es. \texttt{utenti.json}, \texttt{ordini.json}, \texttt{prodotti.json},
\texttt{inventario.json}).

\textbf{Gestione}: la classe \texttt{GsonConfig} fornisce un serializzatore
JSON coerente (helper per \texttt{LocalDate}/\texttt{LocalDateTime}) usato da
tutti i DAO FS per garantire la lettura/scrittura corretta dei record.

\textbf{Configurazione}: la modalità è attivata impostando
\texttt{persistence\_type = FS} nel file di configurazione.

\section{Layer di persistenza DEMO (In-Memory)}
Il layer di persistenza DEMO mantiene i dati esclusivamente nella memoria
volatile (RAM), senza storage permanente. Le strutture sono popolate da valori
predefiniti all'avvio e cancellate al termine dell'esecuzione.

\textbf{Casi d'uso}: è pensato per il testing unitario, lo sviluppo rapido e
le dimostrazioni che non richiedono configurazione di database.

\textbf{Configurazione}: la modalità è attivata impostando
\texttt{persistence\_type = DEMO} (default).
